%\documentclass[DiscStateSpaceToolkit.tex]{subfiles} 
%\begin{document}

\cite{Hansen1985} is an early paper in the RBC literature. At the time it was important for introducing the extensive margin of labour --- working
or not working --- into models (which it did via a trick of using loteries, so that in the aggregate it was back to just being a continuous choice,
and because utility was now quasi-linear the implied Frisch elasticity is infinity; highly unrealistic). Thanks to computational and modelling advances we can now
model both the extensive and intensive margins of labour supply\footnote{The extensive margin is the decision to work or not work, the intensive margin is, 
having decided to work, how many hours to work.} explicitly, and without the need for (highly unrealistic) lotteries, so the model
of \cite{Hansen1985} is no longer empirically useful. It retains historical interest and while now outdated it also provides us with a chance to look
at how the Toolkit output can be used to generate some of the kinds of model statistics that are standard in the literature.

The value function problem to be solved is,
\begin{equation*}
 V(k,z)=\sup_{k',l} \, \{ log(c)+B l+\beta*E[V(k',z')|z] \}
\end{equation*}
subject to
\begin{eqnarray}
 & c+i= exp(z) k^{\theta} l^{1-\theta} \\
 & k'=(1-\delta)k+i \\
 & z'=\rho z +\epsilon' \quad ,\epsilon \overset{iid}{\sim} N(0,\sigma_{\epsilon}^2)
\end{eqnarray}
where $k$ is physical capital, $i$ is investment, $c$ is consumption, $l$ is aggregate labour supply (for the individual household it represents
the probability of working; see \cite{Hansen1985}), $z$ is a productivity shock that follows an AR(1) process; 
$exp(z) k^{\alpha} l^{1-\theta}$ is the production function, $\delta$ is the depreciation rate.


The following code solves this model, using parallelization on GPU, and draws a graph of the value function; it contains many comments describing what is being done
at each step. This code is a copy of $Hansen1985.m$ which you can run and modify to get a better feel for how
to use the toolkit. Below that is the code that defines the return function, $Hansen1985\_{}ReturnFn.m$.

\textbf{Contents of Hansen1985.m}
\lstinputlisting[language=Matlab]{./Examples/MatlabCodes/Hansen1985.m}

\smallskip

\textbf{Contents of Hansen1985$\_{}$ReturnFn.m}
\lstinputlisting[language=Matlab]{./Examples/MatlabCodes/Hansen1985_ReturnFn.m}

%\end{document}